\شروع‌مسئله[48] \متن‌سیاه{پیشنهاد الگوریتم یا داده‌ساختار مناسب}. در هر قسمت از این سؤال از شما خواسته می‌شود یک الگوریتم یا داده‌ساختار مناسب پیشنهاد دهید. این داده‌ساختار یا الگوریتم می‌تواند از مطالب گفته شده در درس، ترکیبی از آن‌ها و یا خلاقیتی باشد. علاوه بر \زیرخط{کلیّت داده‌ساختار و الگوریتم}، \زیرخط{زمان اجرا} و \زیرخط{دلیل انتخاب خود} را نیز در ۴ الی ۵ خط به طور واضح توضیح دهید. همچنین در صورتی که زمان اجرای الگوریتم شما \زیرخط{سرشکن و یا تصادفی} است، آن را بیان کنید. در این سؤال نیاز به اثبات درستی الگوریتم و اثبات زمان اجرای آن نیست.

ورودی برخی سؤال‌ها یک گراف است، اما برای حل مسئله ممکن است لازم باشد گراف جدیدی بسازید. در این موارد روش ساخت گراف را توضیح دهید.

% صورت اصلی مسئله را در این‌جا بنویسید.
\شروع{زیرمسائل}

\زیرمسئله[8]
%2018 Fall Final - P4
در این مسئله، کلمه دنباله‌ای حداکثر ۱۰ حرفی از حروف الفبای فارسی تعریف می‌شود. همچنین جمله دنباله‌ای از کلمه‌ها تعریف می‌شود. دو جملۀ 
$A = (a_1, \dots, a_k)$
و
$B = (b_1, \dots, b_k)$
یکریخت هستند اگر الگوی تکرار کلمات در آن‌ها یکسان باشد، یعنی برای هر $i$ و $j$،
$a_i = a_j$
باشد اگر و تنها اگر
$b_i = b_j$.
در مثال زیر، جمله‌های آ و پ یکریخت هستند زیرا کلمات ۵اُم و ۷اُم در هر کدام از آن‌ها یکسان هستند و بقیه کلماتشان متمایز هستند؛ ولی جمله‌های آ و ب یکریخت نیستند.

آ = 
(''رنج``، ''گرانم``، ''را``، ''به``، ''صحرا``، ''می‌دهم``، ''صحرا``، ''نمی‌گیرد‍``‍‍‍)

ب =
(''رمز``، ''میستان``، ''همه``، ''تو``، ''راز``، ''نیستان``، ''همه``، ''تو``)

پ =
(''اشک``، ''روانم``، ''را``، ''به``، ''دریا``، ''می‌دهم``، ''دریا``، ''نمی‌گیرد``)

لیستی از $n$ جمله به شما داده شده که هر جمله دقیقاً از $k$ کلمه تشکیل شده است. یک الگوریتم با زمان $O(nk)$ طراحی کنید که جمله‌ها را به دسته‌های یکریخت افراز کند.

\پاسخ{
	ایده: مرحلۀ اول تبدیل هر جمله به یک دنباله از اعداد که جمله‌های یکریخت به دنباله‌های یکسان تبدیل شوند و جمله‌های غیریکریخت به دنبال‌های متفاوت. مرحلۀ دوم دسته‌بندی جمله‌ها بر اساس دنبالۀ متناظر به دست آمده در مرحلۀ اول. مرحلۀ اول با جدول درهم‌سازی یا ترای. مرحلۀ دوم با 
	
	ما ابتدا نگاشتی از یک جمله به یک لیست از اعداد پیدا می‌کنیم که دو جمله اگر و تنها اگر یکریخت هستند که لیست مربوط به آن‌ها یکی باشد. برای این کار کلمات جمله را به ترتیب در یک جدول درهم‌سازی (یا ترای) جست‌وجو می‌کنیم. اگر وجود داشت یعنی این کلمه تکراری است و عدد مربوط به آن را در لیست قرار می‌دهیم. در غیر این صورت این کلمه جدید است و آن را در جدول (یا ترای) وارد کرده و عدد جدیدی به آن اختصاص می‌دهیم. اعدادی که اختصاص می‌دهیم اعداد طبیعی به ترتیب با شروع از یک هستند. در این صورت نگاشت یک جمله به یک لیست $O(k)$ زمان می‌گیرد که در صورت استفاده از جدول درهم‌سازی زمان تصادفی است و در صورت استفاده از ترای زمان بدترین حالت است. بنابراین نگاشت همۀ جمله‌ها $O(nk)$ طول می‌کشد. حال کافی است جمله‌ها را توجه به لیست‌شان به دسته‌های هم‌ارزی افراز کنیم. برای این کار می‌توان از یک جدول درهم‌سازی جدید استفاده کرد و وقتی لیست هر جمله را در آن جدول وارد می‌کنیم، اگر تا کنون در آن جدول وجود نداشت یک آرایۀ پویا برای جمله‌های یکریخت با آن با درج اندیس جمله در ابتدای آرایه می‌سازیم و اگر در جدول قبلاً درج شده بود، آرایۀ مربوط به جمله‌های یکریخت با آن را پیدا کردن و اندیس جمله را به انتهای آن آرایه اضافه می‌کنیم.
	
درهم‌سازی اول را می‌توان با ترای جایگزین کرد و درهم‌سازی دوم را با مرتب‌سازی مبنایی.
	
	توجه کنید نیاز به توضیح چگونگی کارکرد جدول درهم‌سازی نیست (مگر این که بخواهیم تغییری روی آن ایجاد کنیم).
	
	همچنین توجه کنید که پیدا کردن دسته‌های هم‌ارزی با زنجیر جدول درهم‌سازی ممکن نیست، زیرا اگر یک کلید را دوبار در جدول وارد کنیم، اولی پاک می‌شود.
	
	زمان درهم‌سازی تصادفی است و ذکر این نکته ضروری است.
 
 



}

\زیرمسئله[8]
%2018 Fall Final - P7
در راستای یک خیابان شرقی-غربی، برج‌های بلندی در حال ساخت و خراب شدن هستند. یک برج در صورتی می‌تواند طلوع خورشید را ببیند که ارتفاع آن از ارتفاع برج‌های سمت شرقش اکیداً بیشتر باشد. داده‌ساختاری برای نگه‌داری برج‌های این خیابان طراحی کنید که اعمال زیر را در زمان $O(\log n)$ پشتیبانی کند ($n$ تعداد برج‌ها در هر لحظه است).
\شروع{فقرات}
\فقره
\متن‌لاتین{\متن‌تایپ{add(adr, hei)}}: یک برج جدید با ارتفاع $h$ که یک عدد صحیح مثبت است، در آدرس $adr$ که آن نیز یک عدد صحیح مثبت است، اضافه شود.
\فقره
\متن‌لاتین{\متن‌تایپ{remove(adr)}}: برجی که در آدرس $adr$ هست از داده‌ساختار حذف شود.
\فقره
\متن‌لاتین{\متن‌تایپ{sunrise(adr)}}:
آیا برجی که در آدرس $adr$ هست طلوع خورشید را می‌بیند؟ اگر بله مقدار ۱ وگرنه مقدار ۰ برگردانده شود.
\پایان{فقرات}

\پاسخ{
	برای این کار از یک درخت AVL افزوده استفاده می‌کنیم که هر رأس بیشینۀ زیردرخت خود را دارد. در این داده‌ساختار همان‌طور که در کلاس گفته شده، می‌توان یک عنصر را در $O(\log n)$ درج یا حذف کرد و یا بیشینۀ یک زیربازه را به دست آورد. توضیح کارکرد و درستی داده‌ساختار در جواب سؤال نیاز نیست.
	
	کافی است برج‌ها را بر حسب $adr$ در درخت درج کنیم و بر اساس ارتفاع برج‌ها اطلاعات افزوده را نگه داریم. در این صورت دو عمل اول به راحتی انجام می‌شود و برای عمل سوم کافی است بیشینه ارتفاع برج‌های قبل از $adr$ را با این داده‌ساختار به دست آوریم و آن را با ارتفاع برجی که در $adr$ قرار دارد مقایسه کنیم. در این صورت هر سه عمل در $O(\log n)$ انجام می‌شوند.
	
	اگر به جای نگه‌داری بیشینۀ زیردرخت در هر رأس، مثلاً بخواهیم بیشترین عنصر در سمت شرق را نگه داریم در موقع به‌روزرسانی به مشکل خواهیم خورد. همچنین ممکن است یکی به اشتباه فرض کنید که رأس‌های سمت شرق در زیردرخت یک رأس وجود دارند که راه‌حل را غلط می‌کند.

زمان درخت پاره‌خطی به حداکثر مقدار $adr$ وابسته است و بنابراین ۶ نمره به آن تعلق می‌گیرد.
}

\زیرمسئله[8]
%2018 Fall Quiz 2 - P7
آقای ووپی در جایی زندگی می‌کند که خیابان‌ها و کوچه‌ها شیب زیادی دارند. او می‌خواهد صبح از خانه تا محل کار را پیاده‌روی کند. محل زندگی او به صورت $n$ تقاطع و $e$ راه یک‌طرفه بین تقاطع‌ها مدل شده است که هر تقاطع یک ارتفاع منحصربه‌فرد دارد. خانه و محل‌کار آقای ووپی هم هر کدام یک تقاطع هستند. مسیر آقای ووپی از خانه به محل کار باید به صورتی باشد که ابتدا ارتفاع مسیر طی شده صعودی باشد، سپس در یک تقاطع کمی استراحت کند و آب بنوشد و سپس ارتفاع ادامۀ مسیر طی شده نزولی باشد. در بین این مسیرها آقای ووپی می‌خواهد کوتاه‌ترین مسیر را پیدا کند. برای این کار یک الگوریتم کارا پیشنهاد دهید.

\پاسخ{
	گراف $G$ را به این صورت می‌سازیم که برای هر تقاطع $x$ دو رأس $x_1$ و $x_2$ قرار می‌دهیم. برای هر راه بین دو تقاطع $x$ و $y$ با وزن $w$، اگر ارتفاع $x$ کمتر باشد، یال جهت‌داری از $x_1$ به $y_1$ با وزن $w$ و اگر ارتفاع $y$ کمتر باشد، یال جهت‌داری از $x_2$ به $y_2$ با وزن $w$ قرار می‌دهیم. همچنین برای هر تقاطع $x$ یال جهت‌داری با وزن صفر از $x_1$ به $x_2$ قرار می‌دهیم. در این گراف رأس‌های با اندیس ۱ مربوط به مسیر صعودی (قبل از توقف) است، یال‌های بین اندیس ۱ و ۲ برای توقف و آب خوردن است و در نهایت رأس‌های با اندیس ۲ مربوط به مسیر نزولی (بعد از توقف) است. توجه کنید گراف جدید بدون دور است (گراف اصلی می‌تواند دور داشته باشد). در این گراف با الگوریتم مربوط به \متن‌لاتین{DAG}، کوتاهترین مسیر را پیدا می‌کنیم. زمان اجرای الگوریتم $O(n+e)$ است.
	
	اگر به جای الگوریتم کوتاه‌ترین مسیر در \متن‌لاتین{DAG}، از دو بار اجرای دکسترا (یک بار از خانه و یک بار برعکس از محل کار) استفاده شود، ۶ نمره داده خواهد شد. همچنین اگر گراف مشابهی ساخته شود ولی در نهایت از دکسترا استفاده شود.
	
	در صورتی که از الگوریتم بلمن‌فورد به صورت درستی استفاده شود یا دکسترا از $n$ رأس 
	دکسترا زده شود ۳ نمره داده خواهد شد.
	
}

\زیرمسئله[8]
رئیس آقای ووپی، آقای استنلی برای پیاده‌روی از منزل به محل کار به دنبال راه ایده‌آل از نظر خودش است. مدل تقاطع‌ها و راه‌ها مثل قسمت قبل است (بدون درنظر گرفتن ارتفاع‌ها). در هر کدام از $e$ راه گفته شده در سؤال قبل، حداقل یک سگ ولگرد زندگی می‌کند. آقای استنلی از سگ‌ها می‌ترسد و برای همین به دنبال راهی هست که با کمترین سگ در مسیر خود روبه‌رو شود. اگر تعداد سگ‌ها $k$ باشد، الگوریتمی با زمان $O(k)$ برای این کار طراحی کنید.

\پاسخ{
	راه‌اصلی: BFS روی گرافی که یال‌ها به خاطر سگ‌ها چند یال شده‌اند.
	
	گراف $G$ را به این صورت می‌سازیم که برای هر تقاطع یک رأس می‌گذاریم و برای هر راه، دو تقاطع دو سر راه را با یک مسیر $x$ یالی به هم وصل می‌کنیم که $x$ تعداد سگ‌های موجود در آن راه است. در این صورت تعداد رأس‌های گراف $G$ برابر با $n+(k-e)$ و تعداد یال‌های این گراف برابر با $k$ خواهد بود. با توجه به این‌که هر راه حداقل یک سگ دارد و به هر تقاطع حداقل یک راه وصل است (وگرنه اسمش تقاطع نبود!)، می‌توان گفت تعداد رأس‌ها و یال‌ها $O(k)$ است و گراف بدون وزن است. حال از جست‌وجوی سطح اول (BFS) استفاده می‌کنیم که مسیر با کمترین یال را از خانه به محل کار آقای استنلی در زمان $O(k)$ پیدا کنیم.
	
توجه کنید گراف برخلاف مسئلۀ قبل دیگر DAG نیست. همچنین استفاده از دکسترا کندتر از BFS خواهد بود (مگر این که به جای فیبوناچی هیپ، از این نکته استفاده کنیم که $d$ها یک عدد صحیح بین صفر تا حداکثر $k$ هستند و از یک جدول دسترسی مستقیم به عنوان صف اولویت استفاده کنیم که تحلیل آسانی هم ندارد. این راه را به احتمال زیاد هیچ کسی ننوشته است).

}

\زیرمسئله[8]
%2013 Quiz 2 P2
%Dijkstra
آقای ووپی می‌خواهد از محل کار خود به یک جلسۀ مهم برود. با توجه به این که ساعت اوج ترافیک است، او می‌خواهد خود را با ماشین به یک ایستگاه مترو برساند و با مترو به محل جلسه برود. مدل تقاطع‌ها و راه‌ها مثل دو قسمت قبل است و $m$ ایستگاه مترو نیز در تقاطع‌های مربوط به خود وجود دارند. هدف این است که آقای ووپی را در سریع‌ترین زمان ممکن به یک ایستگاه مترو برسانیم. الگوریتم بهینه‌ای برای این کار پیشنهاد دهید.

\پاسخ{
	همه مقصدها را به یک رأس جدید وصل می‌کنیم و دکسترا می‌زنیم.
	
}

\زیرمسئله[8]
یک گراف جهت‌دار $G$ داریم که یال‌های آن وزن مثبت یا منفی دارند. الگوریتمی با زمان $O(ne)$ پیشنهاد دهید که مشخص کند آیا مسیری از $s$ به $t$ وجود دارد که وزن آن حداقل $w$ باشد یا خیر؟

\پاسخ{
	بلمن‌فورد روی گرافی که وزن یال‌هایش منفی هستند. یا طولانی‌ترین مسیر یک مسیر ساده است و وزن‌اش بیشتر از $w$ است و یا می‌توان از دور مثبت استفاده کرد. توضیح در کلاس داده شده است.
	
Find a longest path from s to t using a similar algorithm to Problem Set 8 Question 3. Negate all edge weights and run Bellman-Ford from s in O(|V ||E|) time. If a negative weight cycle is not detected, return whether the minimum weight path to t is less than or equal to −w. Otherwise, if a negative weight cycle exists, there exists a path from s to t with weight greater than w: traverse any path from s to t that includes a vertex on the cycle (which exists because the graph is strongly connected), and then splice in as many trips around the cycle as necessary to make the path weight greater than w.
Rubric:
• 12 points for a correct O(|V ||E|) time algorithm • 3 points for running time analysis
• Partial credit may be awarded
	
}

%\زیرمسئله[7]
%گراف وزن‌دار $G$ و رأس $s$ از آن داده شده‌اند. وزن یال‌ها در این گراف می‌تواند منفی یا مثبت باشد. این گراف دقیقاً یک دور مثبت دارد (ولی ممکن است چندین دور منفی داشته باشد). الگوریتمی با زمان $O(n^3)$ پیشنهاد دهید که تمام رئوسی که طولانی‌ترین مسیر از آن‌ها کمتر از $n$ رأس دارد را خروجی دهد.
%
%\پاسخ{
%All simple longest paths will traverse fewer than |V| edges, whereas longest paths which are non-simple will contain a positive weight cycle. Negate edge weights in G and run Bellman-Ford from s to determine whether the positive weight cycle is reachable from s. If not, all vertices reachable from s satisfy the condition as all longest paths traverse fewer than |V | edges. Otherwise, Bellman-Ford finds all vertices reachable from s as well as finding the positive weight cycle reachable from s. Any vertex reachable from arbitrary vertex v on the cycle will not have a longest path from s traversing fewer than |V | edges, so run BFS or DFS to find the set of vertices reachable from v. Then, return the set of vertices reachable from s but not reachable from v. Bellman-Ford runs in O(|V ||E|) = O(|V |3) time and traversal takes O(|V | + |E|) = O(|V |2) time, so this algorithm runs in O(|V |3) time.
%Rubric:
%• 9 points for a correct algorithm. Common approach:
%– 3 points for algorithm to find positive weight cycle – 3 points for finding vertices reachable from cycle – 3 points for returning correct set
%• 2 points for correct running time analysis • 4 points if algorithm is O(|V |3)
%• Partial credit may be awarded
%
%	
%}


%\زیرمسئله[6]
%
%
%
%\پاسخ{
%
%
%	
%}
%
%\زیرمسئله[6]
%یک مجموعه از اعداد متمایز داریم و به دنبال سه تا از این اعداد هستیم که جمع آن‌ها برابر با صفر شود. الگوریتمی با زمان اجرای $O(n^2)$ طراحی کنید که یک سه‌تایی با این خاصیت (در صورت وجود) پیدا کند. 
%به طراحی الگوریتمی با زمان اجرای $O(n^2\log n)$ نیمی از نمره داده خواهد شد. (توجه کنید شما مجاز به استفاده از درهم‌سازی نیستید)
%
%
%\پاسخ{
%راه اول با درهم‌سازی:
%
%
%راه‌دوم: عناصر را در $O(n\log n)$ مرتب می‌کنیم. همۀ عناصر را به عنوان عنصر اول سه‌تایی امتحان می‌کنیم. سپس یک اشاره‌گر مربوط به عنصر دوم به کوچکترین عنصر و یک اشاره‌گر مربوط به عنصر سوم به بزرگترین عنصر قرار می‌دهیم. در هر مرحله اگر جمع این سه عنصر بزرگ‌تر از صفر بود، اشاره‌گر عنصر سوم را حرکت می‌دهیم و اگر بزرگ‌تر از صفر بود، اشاره‌گر عنصر دوم را حرکت می‌دهیم. همچنین به تکراری نبودن عناصر سه‌تایی توجه می‌کنیم. زمان اجرای این الگوریتم برای یک عنصر اول $O(n)$ و در مجموع $O(n^2)$ می‌شود.
%
%}

\پایان{زیرمسائل}